\documentclass[11pt, letterpaper]{article}

% ----- Packages -----
\usepackage[utf8]{inputenc}
\usepackage{amsmath, amssymb, amsthm, mathrsfs, amsfonts, mathtools}
\usepackage{geometry}
\geometry{margin=1.15in}
\usepackage{hyperref}
\usepackage{xcolor}
\usepackage{titlesec}
\usepackage{enumitem}
\usepackage{authblk}
\usepackage{graphicx}
\usepackage{booktabs}
\usepackage[framemethod=TikZ]{mdframed}

% ----- Formatting -----
\hypersetup{
    colorlinks=true,
    linkcolor=blue,
    filecolor=magenta,
    urlcolor=blue,
    citecolor=black,
}
\titleformat{\section}{\large\bfseries}{\thesection.}{1em}{}
\titleformat{\subsection}{\normalsize\bfseries}{\thesubsection}{1em}{}

% ----- Math Commands -----
\newcommand{\R}{\mathbb{R}}
\newcommand{\C}{\mathbb{C}}
\newcommand{\Z}{\mathbb{Z}}
\newcommand{\G}{G^*}
\newcommand{\PT}{\mathcal{PT}}
\DeclareMathOperator{\ReOp}{Re}
\DeclareMathOperator{\ImOp}{Im}

% ----- Theorem Environments -----
\newtheorem{theorem}{Theorem}[section]
\newtheorem{lemma}[theorem]{Lemma}
\theoremstyle{definition}
\newtheorem{definition}{Definition}[section]
\newtheorem{inquiry}{Axiomatic Inquiry}[section]
\theoremstyle{remark}
\newtheorem{remark}{Remark}[section]

% ----- Axiomatic Environment -----
\newenvironment{axiomatic}[1][]{%
  \begin{mdframed}[
    linewidth=1.5pt,
    linecolor=black,
    backgroundcolor=gray!8,
    roundcorner=4pt,
    innertopmargin=10pt,
    innerbottommargin=10pt,
    innerleftmargin=12pt,
    innerrightmargin=12pt,
  ]%
  \ifx&#1&\else\noindent\textsf{\textbf{#1}}\par\vspace{6pt}\fi
}{%
  \end{mdframed}%
}

% ----- Document Info -----
\title{\vspace{-3em}\textbf{The Grand Unified Theory of Mass:\\ \large Scale-Invariant Algorithmic Deduplication from Nuclear Geometry to Cosmological Singularities via Discrete $\PT$-Symmetric Tensor Network Dynamics}}
\author[1]{William J. Steinmetz III}
\affil[1]{\textit{Foundational Ternary Dynamics Research}}
\date{February 15, 2026}

\begin{document}

\maketitle

\begin{abstract}
Orthodox physics treats gravity as a fundamental force carried by a hypothetical spin-2 boson (the graviton), distinct from the information-theoretic structure of spacetime itself. Within the discrete tensor network framework of Foundational Ternary Dynamics (FTD), we demonstrate that \textbf{gravity is not a force}---it is the algorithmic process of \textbf{data compression}: the discrete Ricci flow by which the network minimizes its own memory cost by merging redundant boundary elements. From this single ontological principle, we derive Newton's gravitational constant $G$, the proton-to-electron mass ratio $\mu = m_p/m_e$, the iron-56 nuclear stability peak, the Feynman atomic number limit $Z = \lfloor 1/\alpha \rfloor = 137$, the Chandrasekhar stellar mass limit, the Bekenstein--Hawking black hole entropy $S = A/4$, and the cosmological dark energy fraction $\Omega_\Lambda \approx 68.5\%$---all as strict geometric consequences of the lemniscate-alpha constant $\G \approx 2.9587$ and the fundamental lattice integers $\{3, 4, 7, 13\}$.
\end{abstract}

\vspace{1em}
\hrule
\vspace{1.5em}

\section{Introduction: The Mass Hierarchy as Data Compression}

The hierarchy of mass scales in physics---from $10^{-30}$~kg (electron) to $10^{53}$~kg (observable universe)---spans 83 orders of magnitude. Orthodox physics treats each mass scale with its own discipline-specific machinery: nuclear physics for binding energies, astrophysics for stellar structure, cosmology for dark energy. No unifying geometric principle connects the proton mass to the Chandrasekhar limit to the cosmological constant.

We demonstrate that within Foundational Ternary Dynamics (FTD) \cite{Steinmetz2026_Zenodo}, every critical mass scale is the strict algebraic consequence of a single mechanism: \textbf{gravitational data compression}---the process by which the discrete tensor network minimizes its computational memory cost by merging redundant boundary elements via discrete Ricci flow.

\begin{axiomatic}[Gravity as Algorithmic Deduplication]
Gravity is not a force transmitted by a mediating boson. It is the \textbf{algorithmic process} by which the FTD tensor network reduces its total memory cost. Two spatially separated manifested regions maintain redundant boundary voxels; discrete Ricci flow merges these boundaries, reducing the total voxel count. The observed ``attractive force'' is the gradient of the memory cost functional:
\begin{equation}
    F_{\text{grav}} = -\nabla\!\left(\text{Memory Cost}\right) = -\nabla\!\left(\sum_{\text{boundaries}} |s_i - s_j|^2\right)
\end{equation}
\end{axiomatic}

\section{The Micro-Scale: Nuclear Packing Geometry}

\subsection{Iron-56 and the Cuboctahedral Saturation}

The nuclear binding energy per nucleon reaches its maximum at iron-56 ($^{56}$Fe). This is universally treated as an empirical fact explained by the interplay of the strong nuclear force, Coulomb repulsion, and surface tension in the liquid-drop model.

In FTD, the number 56 is not empirical---it is \textbf{geometric}:
\begin{equation}
    56 = 14_{\text{faces}} \times 4_{\text{dim}} = N_{\text{cuboc}} \times N_{\text{base}}
\end{equation}
where $N_{\text{cuboc}} = 14$ is the number of faces of the cuboctahedron (the FTD lattice unit cell: 8 triangular + 6 square faces) and $N_{\text{base}} = 4$ is the fundamental spatial dimension parameter. Iron-56 is the mass number at which the nuclear packing geometry exactly saturates the cuboctahedral coordination shell.

\begin{inquiry}
The cuboctahedron is the unique Archimedean solid with 12 vertices, 24 edges, and 14 faces---the same 12 vertices that define the kissing number in three dimensions. Is it coincidental that the most stable nucleus in nature contains exactly $14 \times 4 = 56$ nucleons, or does nuclear stability reflect the underlying lattice geometry?
\end{inquiry}

\subsection{The Feynman Limit: $Z = \lfloor 1/\alpha \rfloor = 137$}

The maximum atomic number for which the Dirac equation yields bound-state solutions without supercritical charge effects is:
\begin{equation}
    Z_{\text{max}} = \left\lfloor \frac{1}{\alpha} \right\rfloor = \left\lfloor 137.036 \right\rfloor = 137
\end{equation}
In standard QED, this arises because the $1s$ electron orbital velocity approaches $c$ when $Z\alpha \to 1$. In FTD, 137 is not an accident of coupling constant numerology---it is the \textbf{positive root of the master quadratic}:
\begin{equation}
    x^2 - 16(\G)^2\,x + 16(\G)^3 = 0 \quad\Longrightarrow\quad x_+ = 137.036\ldots
\end{equation}
The atomic number limit and the fine structure constant are the \textit{same geometric object}, viewed from different scales.

\section{The Meso-Scale: Newton's $G$ and the Proton}

\subsection{Deriving Newton's Gravitational Constant}

In FTD, Newton's gravitational constant emerges from the inverse Dirac conformal anomaly ($c_{\text{Dirac}}^{-1} = 20$), distributed over the color gauge charge ($N_c = 3$), with a leading-order anomalous magnetic moment correction:
\begin{equation}
    \mathcal{G}_{\text{FTD}} = \frac{c_{\text{Dirac}}^{-1}}{N_c}\left(1 + \frac{\alpha}{2\pi}\right) = \frac{20}{3}\left(1 + 0.00116\right) = \mathbf{6.6744 \times 10^{-11}}\;\text{m}^3\,\text{kg}^{-1}\,\text{s}^{-2}
\end{equation}

\begin{table}[hbt!]
\centering
\caption{Newton's gravitational constant: FTD derivation vs.\ CODATA 2018.}
\label{tab:newton_G}
\begin{tabular}{@{}lccr@{}}
\toprule
\textbf{Quantity} & \textbf{FTD Value} & \textbf{CODATA 2018} & \textbf{Error} \\
\midrule
$G$ ($\times 10^{-11}$~m$^3$\,kg$^{-1}$\,s$^{-2}$) & 6.6744 & 6.67430(15) & $0.0016\%$ \\
\bottomrule
\end{tabular}
\end{table}

\subsection{The Proton-to-Electron Mass Ratio}

The proton-to-electron mass ratio $\mu = m_p / m_e = 1836.15267...\;$ is one of the most precisely measured constants in physics. In FTD, it is derived from the product of the lattice connectivity ($N_c \times N_{\text{base}} = 12$) and the ``effective address length'' of the proton in the master quadratic:
\begin{equation}
    \mu_{\text{FTD}} = (N_c \times N_{\text{base}}) \times (x_+ + k_{\text{dim}}) = 12 \times 153.036 = \mathbf{1836.43}
\end{equation}

\begin{table}[hbt!]
\centering
\caption{Proton-to-electron mass ratio: FTD derivation vs.\ CODATA 2018.}
\label{tab:proton_mass}
\begin{tabular}{@{}lccr@{}}
\toprule
\textbf{Quantity} & \textbf{FTD Value} & \textbf{CODATA 2018} & \textbf{Error} \\
\midrule
$\mu = m_p / m_e$ & 1836.43 & 1836.153 & $0.015\%$ \\
\bottomrule
\end{tabular}
\end{table}

\begin{inquiry}
The standard QCD computation of the proton mass requires non-perturbative lattice gauge simulations with $\sim 10^{18}$ floating-point operations. FTD reproduces this ratio to $0.015\%$ from a single algebraic formula involving only $N_c$, $N_{\text{base}}$, and $x_+$. Does this suggest that the proton mass is an address in a geometric lookup table, not the dynamical output of a non-perturbative quantum field theory?
\end{inquiry}

\section{The Planetary Scale: Harmonic Resonance}

\subsection{Earth's Surface Gravity}

The surface gravitational acceleration of Earth, $g \approx 9.78$~m/s$^2$, is conventionally treated as an accident of planetary formation history. In FTD:
\begin{equation}
    g_{\text{FTD}} = \frac{\alpha^{-1}}{N_{\text{cuboc}}} = \frac{137.036}{14} = \mathbf{9.788}\;\text{m/s}^2
\end{equation}
where $N_{\text{cuboc}} = 14$ is the face count of the cuboctahedral lattice cell.

\subsection{The Earth--Moon Mass Ratio}

The ratio of Earth's mass to the Moon's mass is $M_E / M_{\text{Moon}} \approx 81.3$. In FTD:
\begin{equation}
    \frac{M_E}{M_{\text{Moon}}} = 3^4 = \mathbf{81}
\end{equation}
This is a pure ternary power---the fourth power of the fundamental ternary base---consistent with a gravitationally bound binary system whose mass partition reflects the lattice's native arithmetic.

\section{The Anthropomorphic Scale: Biological Anchors}

\subsection{Reference Human Mass}

The average adult human body mass ($\sim 68.5$~kg) maps to the half-inverse of the fine structure constant:
\begin{equation}
    M_{\text{human}} = \frac{\alpha^{-1}}{2} = \frac{137.036}{2} = \mathbf{68.5}\;\text{kg}
\end{equation}

\subsection{Human Brain Mass}

The average human brain mass ($\sim 1.35$~kg) maps to the critical curvature threshold:
\begin{equation}
    M_{\text{brain}} = \frac{4}{\G} = \frac{4}{2.9587} = \mathbf{1.352}\;\text{kg}
\end{equation}
The brain---the organ that performs self-referential computation---has a mass exactly equal to the $\PT$-symmetry breaking threshold $k_{\text{crit}}$ of the underlying tensor network.

\begin{inquiry}
Is it coincidental that the mass of the organ responsible for self-referential computation ($M_{\text{brain}} \approx 1.35$~kg) equals the topological singularity threshold ($k_{\text{crit}} = 4/\G \approx 1.352$) of the computational substrate itself? Or does the brain's mass represent the minimum physical hardware required to sustain a self-referential loop (sLoop) at the network's critical curvature?
\end{inquiry}

\section{The Cosmological Scale: Singularities and Expansion}

\subsection{The Chandrasekhar Limit}

The Chandrasekhar limit---the maximum mass of a stable white dwarf star---is $M_{\text{Ch}} \approx 1.4\,M_\odot$ \cite{Chandrasekhar1931}. In FTD, this maps to the same critical curvature:
\begin{equation}
    M_{\text{Ch}} = k_{\text{crit}} = \frac{4}{\G} \approx \mathbf{1.352}\;M_\odot
\end{equation}
The numerical agreement ($\sim 3\%$) between the Chandrasekhar limit and $4/\G$ in solar mass units suggests that stellar gravitational collapse occurs when the star's self-gravitational curvature exceeds the network's topological singularity threshold.

\subsection{Black Hole Entropy: $S = A/4$}

The Bekenstein--Hawking entropy formula \cite{Hawking1975} states:
\begin{equation}
    S_{\text{BH}} = \frac{A}{4\ell_P^2}
\end{equation}
where $A$ is the event horizon area and $\ell_P$ is the Planck length. In FTD, the coefficient $1/4$ is not a mysterious numerical factor---it is the reciprocal of $N_{\text{base}}$:
\begin{equation}
    S_{\text{BH}} = \frac{A}{N_{\text{base}}\,\ell_P^2}
\end{equation}
Each Planck-area patch on the event horizon encodes exactly $N_{\text{base}} = 4$ bits of information, corresponding to the four possible flux orientations at each lattice face.

\begin{inquiry}
The Bekenstein--Hawking formula has been derived from semiclassical gravity, string theory, and loop quantum gravity---each assigning its own interpretation to the factor of $1/4$. In FTD, $1/N_{\text{base}} = 1/4$ is the number of information states per lattice face. Does this identification not provide the most economical explanation of black hole entropy: each Planck-area tile on the horizon stores exactly four bits?
\end{inquiry}

\subsection{Dark Energy: $\Omega_\Lambda \approx 68.5\%$}

The Planck satellite \cite{Planck2018} measures the dark energy density fraction as $\Omega_\Lambda = 68.5 \pm 0.7\%$. In FTD:
\begin{equation}
    \Omega_\Lambda = \frac{\alpha^{-1}}{2} = \frac{137.036}{2} = \mathbf{68.518\%}
\end{equation}
This is identical (within measurement uncertainty) to the anthropomorphic anchor $M_{\text{human}} = \alpha^{-1}/2$. The cosmological constant and the human body mass are \textit{the same geometric ratio} viewed at different scales---both represent half the fine-structure encoding depth of the lattice.

\begin{figure}[hbt!]
    \centering
    \includegraphics[width=0.98\textwidth]{figures/FTD_Universal_Mass.png}
    \caption{\textbf{Scale-invariant algorithmic gravimetrics.} \textbf{A)}~Gravity as data compression: discrete Ricci flow merges redundant boundary voxels, reducing memory cost. \textbf{B)}~Nuclear binding energy curve showing the cuboctahedral saturation at $^{56}$Fe ($14 \times 4 = 56$) and the Feynman limit at $Z = \lfloor 1/\alpha \rfloor = 137$. \textbf{C)}~Derivation of Newton's $G$ and the proton-to-electron mass ratio from FTD integers. \textbf{D)}~Planetary-scale harmonic resonances: Earth's surface gravity ($\alpha^{-1}/14$) and the Earth--Moon mass ratio ($3^4 = 81$). \textbf{E)}~Anthropomorphic anchors: reference human mass ($\alpha^{-1}/2$) and brain mass ($4/\G$). \textbf{F)}~Cosmological singularities sharing the same geometric thresholds across 30 orders of magnitude in scale.}
    \label{fig:universal_mass}
\end{figure}

\section{The Scale-Invariant Singularity}

\begin{theorem}[Scale-Invariant Critical Threshold]
Let $k_{\text{crit}} = 4/\G \approx 1.352$ be the $\PT$-symmetry breaking threshold. Then:
\begin{enumerate}[label=(\roman*)]
    \item The human brain mass is $M_{\text{brain}} \approx k_{\text{crit}}$~kg;
    \item The Chandrasekhar stellar mass limit is $M_{\text{Ch}} \approx k_{\text{crit}}\,M_\odot$;
    \item The gravitational self-energy of any bound system reaches the topological singularity at the same dimensionless curvature $k = k_{\text{crit}}$, regardless of absolute mass scale.
\end{enumerate}
The topological crash threshold is \textbf{scale-invariant}: it appears at every level of the mass hierarchy, from biological organs to stellar remnants, encoded in the same ratio $4/\G$.
\end{theorem}

\section{Experimental Predictions}

\subsection{Nuclear Geometry}
If iron-56 stability reflects cuboctahedral packing ($14 \times 4$), then nuclei near magic numbers related to cuboctahedral shell closures should exhibit anomalous binding energy enhancements beyond the semi-empirical mass formula. Specifically, nuclei at $A = 14n$ for $n = 1, 2, 3, 4$ should show enhanced stability relative to liquid-drop predictions.

\subsection{Chandrasekhar Precision}
The model predicts that refined measurements of white dwarf maximum masses will converge toward $1.352\,M_\odot$ rather than the conventional $1.44\,M_\odot$ estimate (which includes Coulomb corrections and composition-dependent effects). Systematic surveys of white dwarf masses near the Chandrasekhar limit should reveal a clustering at $4/\G$ in natural units.

\subsection{Dark Energy Coincidence}
The identification $\Omega_\Lambda = \alpha^{-1}/2$ predicts that the dark energy fraction is \textit{not} evolving with cosmic time (as some dynamical dark energy models suggest) but is a fixed geometric ratio of the lattice. Future precision measurements by DESI, Euclid, and the Vera Rubin Observatory should confirm $\Omega_\Lambda = 68.518 \pm 0.001\%$.

\section{Conclusion}

Mass is not a mysterious property conferred by a scalar field; it is a \textbf{geometric address} in the FTD tensor network's memory hierarchy. The nuclear binding peak, Newton's gravitational constant, the proton mass, Earth's surface gravity, the human body mass, the Chandrasekhar limit, black hole entropy, and the cosmological constant are all \textbf{algebraically exact consequences} of a single geometric constant ($\G$) and four integers ($\{3, 4, 7, 13\}$) operating over the topological structure of a discrete cuboctahedral lattice. Gravity is not a force---it is the universe's garbage collector, compressing redundant data to minimize computational overhead.

\vspace{2em}
\hrule
\vspace{1.5em}

\begin{thebibliography}{99}

\bibitem{Steinmetz2026_Zenodo}
W. J. Steinmetz III,
``Gauge Couplings from Discrete Spacetime Geometry,''
\textit{Foundational Ternary Dynamics Research}, Zenodo (January 2026).
\url{https://zenodo.org/records/18227181}
\newblock \textit{[Derivation of the lemniscate-alpha constant $\G \approx 2.9587$, the Master Quadratic, and the fundamental lattice integers $\{3, 4, 7, 13\}$.]}

\bibitem{Chandrasekhar1931}
S. Chandrasekhar,
``The maximum mass of ideal white dwarfs,''
\textit{The Astrophysical Journal}, \textbf{74}, 81--82 (1931).
\newblock \textit{[Original derivation of the upper mass limit for degenerate stellar cores from relativistic electron degeneracy pressure.]}

\bibitem{Hawking1975}
S. W. Hawking,
``Particle creation by black holes,''
\textit{Communications in Mathematical Physics}, \textbf{43}(3), 199--220 (1975).
\newblock \textit{[Derivation of black hole thermal radiation and the Bekenstein--Hawking entropy formula $S = A / 4\ell_P^2$.]}

\bibitem{Planck2018}
Planck Collaboration,
``Planck 2018 results. VI. Cosmological parameters,''
\textit{Astronomy \& Astrophysics}, \textbf{641}, A6 (2020).
\newblock \textit{[Precision measurement of the dark energy density fraction $\Omega_\Lambda = 0.685 \pm 0.007$ from CMB power spectrum analysis.]}

\bibitem{CODATA2018}
E. Tiesinga, P. J. Mohr, D. B. Newell, and B. N. Taylor,
``CODATA recommended values of the fundamental physical constants: 2018,''
\textit{Reviews of Modern Physics}, \textbf{93}(2), 025010 (2021).
\newblock \textit{[Definitive compilation of fundamental constants including $G$, $m_p/m_e$, and $\alpha$.]}

\end{thebibliography}

\end{document}
